\documentclass[11pt]{article}
% Preámbulo común de los informes E1, E2 y E3, con aspecto de Word: Calibri,
% paleta del tema de Office, bandas y portada.
% Cada informe lo carga con % Preámbulo común de los informes E1, E2 y E3, con aspecto de Word: Calibri,
% paleta del tema de Office, bandas y portada.
% Cada informe lo carga con % Preámbulo común de los informes E1, E2 y E3, con aspecto de Word: Calibri,
% paleta del tema de Office, bandas y portada.
% Cada informe lo carga con \input{../../latex/preambulo.tex}.
%   \portada[Alumnos]{sigla y curso}{tipo de entrega}{título}{autores}{fecha}
\usepackage[spanish,es-tabla]{babel}
\usepackage[utf8]{inputenc}
\usepackage[T1]{fontenc}
% Carlito: mismas medidas que Calibri, la fuente de Word
\usepackage[sfdefault,lining]{carlito}
\usepackage[letterpaper,margin=2.2cm,top=2.8cm,bottom=2.8cm]{geometry}
\usepackage{xcolor}
\usepackage{graphicx}
\usepackage{array}
\usepackage{colortbl}
\usepackage{tabularx}
\usepackage{microtype}  % reparte el justificado en letras y espacios: menos huecos en columnas angostas
\usepackage{ragged2e}
\setlength{\JustifyingParindent}{0pt}  % sin sangría, como el resto del texto
\usepackage{enumitem}
\usepackage{tikz}
\usetikzlibrary{positioning,arrows.meta,calc,shapes}
\usepackage{eso-pic}
\usepackage{titlesec}
\usepackage{listings}
\usepackage[hidelinks]{hyperref}
\graphicspath{{../../assets/}}

% paleta del tema de Office
\definecolor{banda}{HTML}{4F81BD}
\definecolor{titulo}{HTML}{2E75B5}
\definecolor{subtitulo}{HTML}{5B9BD5}
\definecolor{pizarra}{HTML}{273540}
\definecolor{tablaencabezado}{HTML}{DEEAF6}  % el azul claro de los diagramas
\definecolor{fondosuave}{HTML}{DEEAF6}
\definecolor{grisoffice}{HTML}{7F7F7F}
\definecolor{vinotinto}{HTML}{7B1E2B}

% texto agregado o cambiado por el asistente, pendiente de revisión del equipo
\newcommand{\nuevo}[1]{\textcolor{vinotinto}{#1}}

% títulos como los estilos de título de Word
\titleformat{\section}{\large\bfseries\color{subtitulo}}{}{0pt}{}
\titleformat{\subsection}{\normalsize\bfseries\color{subtitulo}}{}{0pt}{}
\titlespacing*{\section}{0pt}{2ex}{1ex}
\titlespacing*{\subsection}{0pt}{1.5ex}{0.6ex}
\setcounter{secnumdepth}{0}

\setlist{itemsep=0.4ex,leftmargin=1.6em}
\setlist[itemize]{label=\textbullet}
\setlength{\parindent}{0pt}
\setlength{\parskip}{8pt plus 2pt}  % espacio después de párrafo de Word

% tablas tipo Word: grilla completa y encabezado con fondo
% celdas justificadas como el resto del texto; Y es la columna flexible de tabularx
\newcolumntype{L}[1]{>{\justifying\arraybackslash}p{#1}}
\newcolumntype{Y}{>{\justifying\arraybackslash}X}
\newcommand{\enc}[1]{\cellcolor{tablaencabezado}#1}
\newcommand{\iniciartabla}{\renewcommand{\arraystretch}{1.3}\arrayrulecolor{black!70}}

% notas al pie con números también dentro de bloques (minipage), como en Word
\renewcommand{\thempfootnote}{\arabic{mpfootnote}}

% pregunta que abre una respuesta
\newcommand{\pregunta}[1]{\textbf{\color{subtitulo}#1}}

% SQL y pseudocódigo
\lstset{basicstyle=\ttfamily\small,columns=fullflexible,keepspaces=true,
  frame=single,rulecolor=\color{grisoffice},backgroundcolor=\color{fondosuave},
  keywordstyle=\color{titulo}\bfseries,language=SQL,
  morekeywords={BEGIN,COMMIT,ROLLBACK},literate={á}{{\'a}}1 {é}{{\'e}}1
  {í}{{\'i}}1 {ó}{{\'o}}1 {ú}{{\'u}}1 {ñ}{{\~n}}1}

% portada: bandas completas; páginas interiores: media banda arriba a la
% izquierda, media banda abajo a la derecha con el número y logo de marca de agua
\pagestyle{empty}
\newif\ifenportada
\AddToShipoutPictureBG{%
  \ifenportada
    \AtPageUpperLeft{\color{banda}\raisebox{-1.2cm}{\rule{\paperwidth}{0.6cm}}}%
    \AtPageLowerLeft{\color{banda}\raisebox{0.6cm}{\rule{\paperwidth}{0.6cm}}}%
  \else
    \AtPageUpperLeft{\color{banda}\raisebox{-1.2cm}{\rule{0.62\paperwidth}{0.6cm}}}%
    \AtPageLowerLeft{%
      \put(\LenToUnit{0.45\paperwidth},\LenToUnit{0.6cm}){\color{banda}\rule{0.55\paperwidth}{0.6cm}}%
      \put(\LenToUnit{0.45\paperwidth},\LenToUnit{0.78cm}){\hspace{0.2cm}\color{white}\small\thepage}%
      \put(\LenToUnit{0.7cm},\LenToUnit{0.5cm}){\logomarcaagua}}%
  \fi}
\newcommand{\logomarcaagua}{\tikz\node[opacity=0.08,inner sep=0]{\includegraphics[width=2cm]{logo_uc.png}};}

% \portada[rótulo]{sigla y curso}{tipo de entrega}{título}{autores}{fecha}
\newcommand{\portada}[6][Alumno]{%
  \global\enportadatrue
  \begin{titlepage}
    \centering
    \vspace*{2.2cm}
    \includegraphics[width=8cm]{logo_uc.png}\par\vspace{1.6cm}
    {\large\bfseries\color{titulo} Magíster en Inteligencia Artificial - MIA\par}\vspace{0.4cm}
    {\large\bfseries\color{subtitulo} Instituto de Ingeniería Matemática y Computacional - IMC\par}\vspace{0.5cm}
    {\LARGE\bfseries\color{pizarra} #2\par}\vspace{0.5cm}
    {\itshape\bfseries\color{subtitulo} #3\par}\vspace{0.2cm}
    {\large #4\par}\vspace{0.2cm}
    {\color{banda}\rule{0.8\textwidth}{0.6pt}\par}
    \vspace{2cm}
    \raggedright
    {\textbf{#1:} #5\par}\vspace{0.3cm}
    {\textbf{Fecha de entrega:} #6\par}
    \vfill
  \end{titlepage}
  \global\enportadafalse}
.
%   \portada[Alumnos]{sigla y curso}{tipo de entrega}{título}{autores}{fecha}
\usepackage[spanish,es-tabla]{babel}
\usepackage[utf8]{inputenc}
\usepackage[T1]{fontenc}
% Carlito: mismas medidas que Calibri, la fuente de Word
\usepackage[sfdefault,lining]{carlito}
\usepackage[letterpaper,margin=2.2cm,top=2.8cm,bottom=2.8cm]{geometry}
\usepackage{xcolor}
\usepackage{graphicx}
\usepackage{array}
\usepackage{colortbl}
\usepackage{tabularx}
\usepackage{microtype}  % reparte el justificado en letras y espacios: menos huecos en columnas angostas
\usepackage{ragged2e}
\setlength{\JustifyingParindent}{0pt}  % sin sangría, como el resto del texto
\usepackage{enumitem}
\usepackage{tikz}
\usetikzlibrary{positioning,arrows.meta,calc,shapes}
\usepackage{eso-pic}
\usepackage{titlesec}
\usepackage{listings}
\usepackage[hidelinks]{hyperref}
\graphicspath{{../../assets/}}

% paleta del tema de Office
\definecolor{banda}{HTML}{4F81BD}
\definecolor{titulo}{HTML}{2E75B5}
\definecolor{subtitulo}{HTML}{5B9BD5}
\definecolor{pizarra}{HTML}{273540}
\definecolor{tablaencabezado}{HTML}{DEEAF6}  % el azul claro de los diagramas
\definecolor{fondosuave}{HTML}{DEEAF6}
\definecolor{grisoffice}{HTML}{7F7F7F}
\definecolor{vinotinto}{HTML}{7B1E2B}

% texto agregado o cambiado por el asistente, pendiente de revisión del equipo
\newcommand{\nuevo}[1]{\textcolor{vinotinto}{#1}}

% títulos como los estilos de título de Word
\titleformat{\section}{\large\bfseries\color{subtitulo}}{}{0pt}{}
\titleformat{\subsection}{\normalsize\bfseries\color{subtitulo}}{}{0pt}{}
\titlespacing*{\section}{0pt}{2ex}{1ex}
\titlespacing*{\subsection}{0pt}{1.5ex}{0.6ex}
\setcounter{secnumdepth}{0}

\setlist{itemsep=0.4ex,leftmargin=1.6em}
\setlist[itemize]{label=\textbullet}
\setlength{\parindent}{0pt}
\setlength{\parskip}{8pt plus 2pt}  % espacio después de párrafo de Word

% tablas tipo Word: grilla completa y encabezado con fondo
% celdas justificadas como el resto del texto; Y es la columna flexible de tabularx
\newcolumntype{L}[1]{>{\justifying\arraybackslash}p{#1}}
\newcolumntype{Y}{>{\justifying\arraybackslash}X}
\newcommand{\enc}[1]{\cellcolor{tablaencabezado}#1}
\newcommand{\iniciartabla}{\renewcommand{\arraystretch}{1.3}\arrayrulecolor{black!70}}

% notas al pie con números también dentro de bloques (minipage), como en Word
\renewcommand{\thempfootnote}{\arabic{mpfootnote}}

% pregunta que abre una respuesta
\newcommand{\pregunta}[1]{\textbf{\color{subtitulo}#1}}

% SQL y pseudocódigo
\lstset{basicstyle=\ttfamily\small,columns=fullflexible,keepspaces=true,
  frame=single,rulecolor=\color{grisoffice},backgroundcolor=\color{fondosuave},
  keywordstyle=\color{titulo}\bfseries,language=SQL,
  morekeywords={BEGIN,COMMIT,ROLLBACK},literate={á}{{\'a}}1 {é}{{\'e}}1
  {í}{{\'i}}1 {ó}{{\'o}}1 {ú}{{\'u}}1 {ñ}{{\~n}}1}

% portada: bandas completas; páginas interiores: media banda arriba a la
% izquierda, media banda abajo a la derecha con el número y logo de marca de agua
\pagestyle{empty}
\newif\ifenportada
\AddToShipoutPictureBG{%
  \ifenportada
    \AtPageUpperLeft{\color{banda}\raisebox{-1.2cm}{\rule{\paperwidth}{0.6cm}}}%
    \AtPageLowerLeft{\color{banda}\raisebox{0.6cm}{\rule{\paperwidth}{0.6cm}}}%
  \else
    \AtPageUpperLeft{\color{banda}\raisebox{-1.2cm}{\rule{0.62\paperwidth}{0.6cm}}}%
    \AtPageLowerLeft{%
      \put(\LenToUnit{0.45\paperwidth},\LenToUnit{0.6cm}){\color{banda}\rule{0.55\paperwidth}{0.6cm}}%
      \put(\LenToUnit{0.45\paperwidth},\LenToUnit{0.78cm}){\hspace{0.2cm}\color{white}\small\thepage}%
      \put(\LenToUnit{0.7cm},\LenToUnit{0.5cm}){\logomarcaagua}}%
  \fi}
\newcommand{\logomarcaagua}{\tikz\node[opacity=0.08,inner sep=0]{\includegraphics[width=2cm]{logo_uc.png}};}

% \portada[rótulo]{sigla y curso}{tipo de entrega}{título}{autores}{fecha}
\newcommand{\portada}[6][Alumno]{%
  \global\enportadatrue
  \begin{titlepage}
    \centering
    \vspace*{2.2cm}
    \includegraphics[width=8cm]{logo_uc.png}\par\vspace{1.6cm}
    {\large\bfseries\color{titulo} Magíster en Inteligencia Artificial - MIA\par}\vspace{0.4cm}
    {\large\bfseries\color{subtitulo} Instituto de Ingeniería Matemática y Computacional - IMC\par}\vspace{0.5cm}
    {\LARGE\bfseries\color{pizarra} #2\par}\vspace{0.5cm}
    {\itshape\bfseries\color{subtitulo} #3\par}\vspace{0.2cm}
    {\large #4\par}\vspace{0.2cm}
    {\color{banda}\rule{0.8\textwidth}{0.6pt}\par}
    \vspace{2cm}
    \raggedright
    {\textbf{#1:} #5\par}\vspace{0.3cm}
    {\textbf{Fecha de entrega:} #6\par}
    \vfill
  \end{titlepage}
  \global\enportadafalse}
.
%   \portada[Alumnos]{sigla y curso}{tipo de entrega}{título}{autores}{fecha}
\usepackage[spanish,es-tabla]{babel}
\usepackage[utf8]{inputenc}
\usepackage[T1]{fontenc}
% Carlito: mismas medidas que Calibri, la fuente de Word
\usepackage[sfdefault,lining]{carlito}
\usepackage[letterpaper,margin=2.2cm,top=2.8cm,bottom=2.8cm]{geometry}
\usepackage{xcolor}
\usepackage{graphicx}
\usepackage{array}
\usepackage{colortbl}
\usepackage{tabularx}
\usepackage{microtype}  % reparte el justificado en letras y espacios: menos huecos en columnas angostas
\usepackage{ragged2e}
\setlength{\JustifyingParindent}{0pt}  % sin sangría, como el resto del texto
\usepackage{enumitem}
\usepackage{tikz}
\usetikzlibrary{positioning,arrows.meta,calc,shapes}
\usepackage{eso-pic}
\usepackage{titlesec}
\usepackage{listings}
\usepackage[hidelinks]{hyperref}
\graphicspath{{../../assets/}}

% paleta del tema de Office
\definecolor{banda}{HTML}{4F81BD}
\definecolor{titulo}{HTML}{2E75B5}
\definecolor{subtitulo}{HTML}{5B9BD5}
\definecolor{pizarra}{HTML}{273540}
\definecolor{tablaencabezado}{HTML}{DEEAF6}  % el azul claro de los diagramas
\definecolor{fondosuave}{HTML}{DEEAF6}
\definecolor{grisoffice}{HTML}{7F7F7F}
\definecolor{vinotinto}{HTML}{7B1E2B}

% texto agregado o cambiado por el asistente, pendiente de revisión del equipo
\newcommand{\nuevo}[1]{\textcolor{vinotinto}{#1}}

% títulos como los estilos de título de Word
\titleformat{\section}{\large\bfseries\color{subtitulo}}{}{0pt}{}
\titleformat{\subsection}{\normalsize\bfseries\color{subtitulo}}{}{0pt}{}
\titlespacing*{\section}{0pt}{2ex}{1ex}
\titlespacing*{\subsection}{0pt}{1.5ex}{0.6ex}
\setcounter{secnumdepth}{0}

\setlist{itemsep=0.4ex,leftmargin=1.6em}
\setlist[itemize]{label=\textbullet}
\setlength{\parindent}{0pt}
\setlength{\parskip}{8pt plus 2pt}  % espacio después de párrafo de Word

% tablas tipo Word: grilla completa y encabezado con fondo
% celdas justificadas como el resto del texto; Y es la columna flexible de tabularx
\newcolumntype{L}[1]{>{\justifying\arraybackslash}p{#1}}
\newcolumntype{Y}{>{\justifying\arraybackslash}X}
\newcommand{\enc}[1]{\cellcolor{tablaencabezado}#1}
\newcommand{\iniciartabla}{\renewcommand{\arraystretch}{1.3}\arrayrulecolor{black!70}}

% notas al pie con números también dentro de bloques (minipage), como en Word
\renewcommand{\thempfootnote}{\arabic{mpfootnote}}

% pregunta que abre una respuesta
\newcommand{\pregunta}[1]{\textbf{\color{subtitulo}#1}}

% SQL y pseudocódigo
\lstset{basicstyle=\ttfamily\small,columns=fullflexible,keepspaces=true,
  frame=single,rulecolor=\color{grisoffice},backgroundcolor=\color{fondosuave},
  keywordstyle=\color{titulo}\bfseries,language=SQL,
  morekeywords={BEGIN,COMMIT,ROLLBACK},literate={á}{{\'a}}1 {é}{{\'e}}1
  {í}{{\'i}}1 {ó}{{\'o}}1 {ú}{{\'u}}1 {ñ}{{\~n}}1}

% portada: bandas completas; páginas interiores: media banda arriba a la
% izquierda, media banda abajo a la derecha con el número y logo de marca de agua
\pagestyle{empty}
\newif\ifenportada
\AddToShipoutPictureBG{%
  \ifenportada
    \AtPageUpperLeft{\color{banda}\raisebox{-1.2cm}{\rule{\paperwidth}{0.6cm}}}%
    \AtPageLowerLeft{\color{banda}\raisebox{0.6cm}{\rule{\paperwidth}{0.6cm}}}%
  \else
    \AtPageUpperLeft{\color{banda}\raisebox{-1.2cm}{\rule{0.62\paperwidth}{0.6cm}}}%
    \AtPageLowerLeft{%
      \put(\LenToUnit{0.45\paperwidth},\LenToUnit{0.6cm}){\color{banda}\rule{0.55\paperwidth}{0.6cm}}%
      \put(\LenToUnit{0.45\paperwidth},\LenToUnit{0.78cm}){\hspace{0.2cm}\color{white}\small\thepage}%
      \put(\LenToUnit{0.7cm},\LenToUnit{0.5cm}){\logomarcaagua}}%
  \fi}
\newcommand{\logomarcaagua}{\tikz\node[opacity=0.08,inner sep=0]{\includegraphics[width=2cm]{logo_uc.png}};}

% \portada[rótulo]{sigla y curso}{tipo de entrega}{título}{autores}{fecha}
\newcommand{\portada}[6][Alumno]{%
  \global\enportadatrue
  \begin{titlepage}
    \centering
    \vspace*{2.2cm}
    \includegraphics[width=8cm]{logo_uc.png}\par\vspace{1.6cm}
    {\large\bfseries\color{titulo} Magíster en Inteligencia Artificial - MIA\par}\vspace{0.4cm}
    {\large\bfseries\color{subtitulo} Instituto de Ingeniería Matemática y Computacional - IMC\par}\vspace{0.5cm}
    {\LARGE\bfseries\color{pizarra} #2\par}\vspace{0.5cm}
    {\itshape\bfseries\color{subtitulo} #3\par}\vspace{0.2cm}
    {\large #4\par}\vspace{0.2cm}
    {\color{banda}\rule{0.8\textwidth}{0.6pt}\par}
    \vspace{2cm}
    \raggedright
    {\textbf{#1:} #5\par}\vspace{0.3cm}
    {\textbf{Fecha de entrega:} #6\par}
    \vfill
  \end{titlepage}
  \global\enportadafalse}

\begin{document}
\portada[Alumnos]{INF3862-1 Plataformas Big Data para Data Science}
{Proyecto del curso, Entrega 2: Ingesta y procesamiento}
{Hábitos de actividad y alimentación asociados a mejores indicadores de salud}
{Miguel Angel Cervantes Buleje y Roger Ruben Reyes Miranda}
{30 de noviembre de 2026}

\begingroup\color{vinotinto}
\section{5.1 Diseño del pipeline}

\pregunta{¿ETL o ELT?} ELT. Los archivos de los relojes se guardan tal como
llegan y se transforman después, en el destino analítico\footnote{Hagedorn,
\emph{Clase 02}, pág. 26, sobre ELT: \textquotedbl{}Conservar datos extraídos y
transformar en el destino\textquotedbl{}.}, por dos condiciones del caso.
Primero, corregir una regla de calidad (por ejemplo, el mínimo de mediciones
por hora de la frecuencia en reposo) o la unificación de los dos modelos de
reloj, que envían los mismos campos con nombres y formatos distintos, exige
recalcular desde los originales. Segundo, el destino puede transformar: DuckDB
procesa bronze y silver en 16.4 s (5.2), de modo que no hace falta un entorno
de transformación previo\footnote{Hagedorn, \emph{Clase 02}, pág. 26, sobre
ELT: \textquotedbl{}El destino debe poder procesar y gobernar los
datos\textquotedbl{}.}.

\pregunta{¿Qué contiene cada capa, en qué formato y quién la consume?}

\begingroup\iniciartabla\small
\begin{tabularx}{\textwidth}{|L{1.5cm}|Y|L{3.4cm}|L{3.6cm}|}\hline
\enc{\textbf{Capa}} & \enc{\textbf{Contenido}} & \enc{\textbf{Formato}} &
\enc{\textbf{Quién la consume}}\\\hline
Bronze & Los archivos del curso sin cambios: mediciones por día y hora de
recepción, dispositivos, sesiones, comidas y recetas & JSON comprimido, CSV,
JSONL y Parquet & La tarea silver y cualquier reproceso\\\hline
Silver & Mediciones unificadas, sin repetidos y tomadas mientras el reloj
estaba vinculado a su usuario;
sesiones y comidas que cumplen las reglas de calidad & Parquet con ZSTD;
mediciones particionadas por fecha de medición & La tarea gold y quien revisa
la calidad\\\hline
Gold & Los tres hechos de la E1 (por usuario y día, por usuario, reloj y día,
y por comida) y las dimensiones de fecha, usuario, reloj y receta & Parquet con ZSTD; hechos particionados por
mes & El dashboard del equipo de nutrición (consultas 4 y 6), el analista
(consulta 5) y la tabla de modelamiento de la E3\\\hline
\end{tabularx}
\endgroup

\pregunta{¿Qué tareas tiene el DAG, cómo dependen entre sí y con qué
frecuencia corre?} Corre una vez al día, después de recibir los archivos
horarios del día anterior (uno por hora con datos), y cada tarea empieza cuando termina la anterior.
Las tareas con fondo azul están implementadas en \texttt{src/pipeline.py}, y el
notebook las ejecuta en orden; las de borde discontinuo son la operación diaria
propuesta.

\begin{figure}[H]
\centering
\begin{tikzpicture}[node distance=0.45cm,
  caja/.style={draw=banda,fill=fondosuave,rounded corners,align=center,
    font=\footnotesize,text width=2.25cm,minimum height=1.25cm},
  propuesta/.style={caja,dashed,fill=white},
  flecha/.style={-{Stealth},draw=pizarra,thick}]
\node[propuesta] (inv) {inventariar\_lote\\archivos del día y conteos};
\node[caja,right=of inv] (bro) {cargar\_bronze\\leer y unificar los dos esquemas};
\node[caja,right=of bro] (sil) {construir\_silver\\reglas de calidad y partición por fecha};
\node[caja,right=of sil] (gol) {construir\_gold\\hechos y dimensiones de la E1};
\node[propuesta,right=of gol] (pub) {validar y publicar\\conteos esperados y consultas 4 a 6};
\draw[flecha] (inv) -- (bro); \draw[flecha] (bro) -- (sil);
\draw[flecha] (sil) -- (gol); \draw[flecha] (gol) -- (pub);
\end{tikzpicture}
\caption{DAG diario del pipeline para el día $d$.}
\end{figure}

\pregunta{Si una carga falla a la mitad y se vuelve a ejecutar, ¿cómo se evita
duplicar registros?} Cada tarea borra y vuelve a escribir completas las
carpetas que produce, y silver deja una sola medición por reloj y marca de
tiempo antes de escribir. Por lo tanto, repetir una tarea después de una falla
deja el mismo estado final\footnote{Hagedorn, \emph{Clase 02}, pág. 24:
\textquotedbl{}Idempotencia: repetir la misma carga produce el mismo estado
final\textquotedbl{}.}. El notebook lo comprueba: al repetir silver y gold, las
siete tablas gold quedan con las mismas filas (17497 en
\texttt{hecho\_usuario\_dia}). Mientras se reescribe una carpeta, el dashboard
podría leerla a medio escribir; por eso la última tarea publica solo después de
validar los conteos. La E1 proponía tablas Delta, que confirman una
reescritura como una unidad; el prototipo usa carpetas de Parquet sin Delta,
porque corre en una máquina y nadie consulta mientras escribe. Este cambio y
los demás respecto de la E1 están en \texttt{decisiones.md}.

\pregunta{¿Cómo se tratan los registros que llegan atrasados?} No se
descartan: silver los guarda en la partición de su fecha de medición, no en la
de su recepción. El prototipo procesa el lote completo en cada ejecución. En
la operación diaria, como el atraso máximo observado es de 4 días (5.2), la
ejecución del día $d$ vuelve a escribir las particiones de los días $d-4$ a $d$, y gold
recalcula los meses que tocan esos días. Una medición que llegue después de
esa ventana queda en bronze y entra en el siguiente reproceso completo, que
se ejecuta cuando cambia una regla de calidad.

\section{5.2 Implementación mínima}

El notebook \texttt{notebooks/proyecto.ipynb} corre de principio a fin con
DuckDB, en el contenedor del repositorio, sobre los datos completos del curso
(61 días de los 300 usuarios con reloj): el lote pesa 296 MB, de los cuales
115.5 MB son mediciones y 178.7 MB recetas (E1), y bronze y silver
tardan 16.4 s, de modo que no hace falta una muestra. El README explica cómo
obtener los datos y ejecutarlo.

\pregunta{¿Qué reglas de calidad se aplican y cuántos registros afecta cada
una?} Seis reglas, sobre las 15540132 mediciones leídas de 1449 archivos:

\begingroup\iniciartabla\small
\begin{tabularx}{\textwidth}{|L{2.9cm}|Y|>{\RaggedRight\arraybackslash}p{3.6cm}|}\hline
\enc{\textbf{Regla}} & \enc{\textbf{Qué hace y qué consecuencia tiene}} &
\enc{\textbf{Registros}}\\\hline
Campos obligatorios & Unifica los dos esquemas de reloj y exige reloj y
marca de tiempo en cada medición & 0\\\hline
Claves que no calzan & Descarta las mediciones de relojes que no están en
el registro de dispositivos o tomadas fuera del período en que el reloj estaba
vinculado a su usuario, y comidas cuya receta
no está en el catálogo, porque no tienen nutrientes & 0 y 0 mediciones; 48
comidas\\\hline
Duplicados & Deja una medición por reloj y marca de tiempo, la primera
recibida: ninguna repetida tiene valores distintos, de modo que son lotes
reenviados. Deja una sesión por \texttt{session\_id}, para no sumar dos veces
sus minutos & 345599 mediciones (2.2\%); 48 sesiones\\\hline
Nulos & Conserva las mediciones sin frecuencia cardíaca, porque sus pasos
valen, y las deja fuera del promedio de reposo & 349931 (2.3\% de las
15194533 sin repetidos)\\\hline
Rangos válidos & Supuesto: frecuencia entre 20 y 250 latidos por minuto;
pasos y calorías no negativos & 0\\\hline
Atrasos & Guarda cada medición en el día en que se midió, no en el que llegó &
8957451 llegan más de 1 hora después de medirse; 2185094, otro día (hasta 4 días)\\\hline
\end{tabularx}
\endgroup

La E1 cuenta 357927 frecuencias nulas porque las cuenta antes de quitar los
repetidos. Silver queda con 15194533 mediciones en 61 particiones por fecha de
medición (52.3 MB, frente a 115.5 MB de JSON comprimido en bronze), 4792
sesiones y 22058 comidas.

\pregunta{¿Cómo se construyen la tabla de hechos y las dimensiones?} Con los
nombres y claves del anexo B de la E1, en Parquet y con los hechos
particionados por mes:
\begin{itemize}
  \item \textbf{\texttt{hecho\_usuario\_dia}:} 17497 filas, una por usuario y
  día local con mediciones o comidas. En los 1167 días con comidas pero sin
  mediciones (674 anteriores a la vinculación del reloj), los pasos quedan
  nulos y no en cero, porque el reloj no midió. La frecuencia en reposo, el
  promedio entre las 2:00 y las 4:59 hora local fuera de sesiones (E1), existe
  en 11823 días de 216 usuarios, los mismos de la E1: sus 11864 días con datos
  menos los 41 sin mediciones de madrugada.
  \item \textbf{\texttt{hecho\_usuario\_reloj\_dia}:} 16330 filas; separa los
  pasos de cada reloj en los días en que un usuario cambia de modelo.
  \item \textbf{\texttt{hecho\_comida}:} 22058 comidas con sus nutrientes
  consumidos (nutrientes por porción por porciones).
  \item \textbf{Dimensiones:} 61 fechas, 300 usuarios, 315 relojes y las 20654
  recetas que aparecen en las comidas.
\end{itemize}

\pregunta{¿Qué responden las tres consultas analíticas de la E1?} Corren sobre
Parquet con el SQL del anexo A de la E1. El único ajuste es que la consulta 6
agrupa también por la clave de la receta, para no juntar recetas con el mismo
nombre: 469 nombres corresponden a dos o más recetas.
\begin{itemize}
  \item \textbf{Consulta 4, hábitos por cuartil:} los 216 usuarios se ordenan
  por frecuencia en reposo en cuartiles de 54, de 51.4 (cuartil 1, el mejor
  indicador) a 69.3 latidos por minuto. Las cifras son promedios diarios. La
  actividad no ordena a los cuartiles: el cuartil 1 camina 8515 pasos diarios
  y el cuartil 4, 8822; el que menos camina es el cuartil 2 (8070). El cuartil
  4 consume más sodio (4588 mg), grasa saturada (53.6 g) y fibra (18.1 g) que
  los otros tres.
  \item \textbf{Consulta 5, pasos semanales por modelo:} los dos modelos reportan pasos comparables. La
  diferencia semanal va de 23 a 458 pasos, y FitBand 2 supera a Pulso S en 7
  de 10 semanas. Por lo tanto, la consulta 4 junta usuarios de ambos modelos
  sin corregir sus pasos.
  \item \textbf{Consulta 6, recetas más preparadas por cuartil:} ninguna receta se repite más de 4 veces en un
  cuartil, entre 15889 comidas de usuarios con cuartil. El ranking de recetas
  no distingue a los cuartiles; el dashboard los compara mejor por nutrientes.
\end{itemize}
Son descripciones de datos sintéticos, no relaciones causales ni
diagnósticos.

\section{5.3 Medir y explicar}

\pregunta{¿Cuánto tarda la consulta 4 sobre CSV y sobre Parquet, y por qué?}
Sobre Parquet tarda 29.01 ms, 10.1 veces menos que sobre CSV (293.42 ms), y
ocupa 0.53 MB, 4.1 veces menos que el CSV (2.15 MB), con el mismo resultado.
Son medianas de 7 ejecuciones en el contenedor (12 CPU y 10.2 GB de memoria).
\begin{itemize}
  \item \textbf{Organización por columnas:} la consulta usa 8 de las 18
  columnas de \texttt{hecho\_usuario\_dia}. DuckDB lee de Parquet solo esas
  8\footnote{Hagedorn, \emph{Clase 01}, pág. 26: \textquotedbl{}Un motor
  columnar puede leer solo los atributos necesarios\textquotedbl{}.}, ya con
  su tipo. El CSV guarda filas de texto: hay que leer las 18 columnas de cada
  fila y convertir el texto en números y fechas.
  \item \textbf{Compresión:} Parquet guarda juntos los valores de cada
  columna, que se parecen entre sí, y los comprime con ZSTD; por eso ocupa 4.1
  veces menos en disco\footnote{Hagedorn, \emph{Clase 01},
  pág. 27: \textquotedbl{}Los valores similares se comprimen
  mejor\textquotedbl{}, con \textquotedbl{}Menos espacio y menos bytes
  leídos\textquotedbl{}.}.
  \item \textbf{Poda de particiones:} no participa en esta medición. La
  consulta pide el período completo y lee los 3 meses; con un filtro de un
  mes leería una de las 3 carpetas\footnote{Hagedorn, \emph{Clase 01}, pág. 40:
  el particionamiento reduce los \textquotedbl{}Datos a examinar, si el filtro
  permite podar\textquotedbl{}.}.
\end{itemize}

\section{5.4 ¿Realmente necesitamos Big Data?}

\pregunta{¿Qué volumen tendría la app con todos los usuarios de Food.com con
reloj, a 1 y 3 años, y en un escenario x100?} Supuesto: los 271907 usuarios
que reseñan en Food.com tienen reloj y generan, como los 300 actuales, 849.2
mediciones por día. El escenario x100 multiplica esos usuarios por 100. Los
tamaños y tiempos se extrapolan en forma lineal desde la corrida del notebook
(7.43 bytes por medición en bronze y 3.43 en silver; 16.4 s de bronze y silver
para 15540132 mediciones): son estimaciones, no mediciones. Bronze y silver
cuentan solo las mediciones, que son las que crecen con los usuarios.

\begingroup\iniciartabla\small
\begin{tabularx}{\textwidth}{|Y|r|r|r|r|r|}\hline
\enc{\textbf{Escenario}} & \enc{\textbf{Mediciones por día}} & \enc{\textbf{Mediciones en total}} &
\enc{\textbf{Bronze}} & \enc{\textbf{Silver}} & \enc{\textbf{Lote diario}}\\\hline
Hoy: 300 usuarios, 61 días & 0.25 millones & 0.02 mil millones & 0.12 GB &
0.05 GB & menos de 1 s\\\hline
Food.com, 1 año & 230.9 millones & 84.28 mil millones & 626.39 GB & 289.41 GB
& 4.1 min\\\hline
Food.com, 3 años & 230.9 millones & 252.84 mil millones & 1879.17 GB & 868.22
GB & 4.1 min\\\hline
Food.com x100, 3 años & 23090 millones & 25283.55 mil millones & 187916.5 GB &
86822.28 GB & 406.1 min\\\hline
\end{tabularx}
\endgroup

\pregunta{¿Seguiría funcionando el pipeline en una máquina? ¿En qué punto, y
por qué recurso, dejaría de alcanzar?} Con todos los usuarios de Food.com, sí,
si procesa cada archivo horario por separado; con el volumen x100, no.
\begin{itemize}
  \item \textbf{Memoria:} es el primer límite del diseño actual. Bronze y
  silver guardan las mediciones en tablas en memoria, y la corrida usó 6.1 GB
  de 10.2 GB con 15540132 mediciones. Un lote diario con todos los usuarios
  tiene 15 veces más mediciones: extrapolando en forma lineal, necesitaría
  unos 91 GB, nueve veces la memoria del contenedor. La corrección es
  procesar cada archivo horario por separado (9.6 millones de mediciones,
  menos que los 61 días actuales) y escribirlo directo en Parquet. En el
  escenario x100, un solo archivo horario tendría 962.1 millones de
  mediciones y necesitaría unos 378 GB: ya no cabe en la memoria de una
  máquina.
  \item \textbf{Disco:} a 3 años, bronze y silver suman 2747.39 GB, de modo
  que hace falta una máquina con varios TB de disco.
  \item \textbf{Tiempo:} el lote diario tardaría 4.1 minutos, y volver a
  procesar los 3 años, unas 74.1 horas. En el escenario x100, el lote diario
  tardaría 406.1 minutos y volver a procesar los 3 años, unas 7411.8 horas
  (309 días).
\end{itemize}

\pregunta{Si se llevara a Spark, ¿qué operaciones provocarían shuffle y cómo se
particionarían los datos para reducirlo?} Agrupar por una clave reúne en un
mismo nodo los datos que quedaron en particiones distintas\footnote{Hagedorn,
\emph{Clase 04}, pág. 15: \textquotedbl{}Agrupar es wide: el shuffle reúne los
aportes de cada país\textquotedbl{}.}:
\begin{itemize}
  \item \textbf{Con shuffle, si los datos no están particionados por reloj:}
  quitar repetidos por reloj y marca de tiempo; sumar pasos y promediar la
  frecuencia cardíaca por usuario, reloj, día y hora; y ordenar a
  los usuarios en cuartiles, que lleva todos los perfiles a una sola
  partición, aunque son solo 216 filas.
  \item \textbf{Sin shuffle, con broadcast:} unir las mediciones con los 315
  relojes (para conocer usuario, modelo y zona horaria) y con las 4792
  sesiones (para excluir de la frecuencia en reposo las mediciones de una
  sesión), y las comidas con las 20654 recetas de
  \texttt{dim\_receta}: copiar la tabla pequeña a cada nodo evita redistribuir
  la grande\footnote{Hagedorn, \emph{Clase 04}, pág. 21: \textquotedbl{}Con
  broadcast, mover la tabla pequeña evita redistribuir la grande para este
  join\textquotedbl{}.}.
  \item \textbf{Partición propuesta en Spark:} por fecha de medición, como
  silver hoy, y además por reloj. Todas las
  copias de una medición repetida quedan en la misma partición, de modo que
  quitar repetidos y sumar por reloj y día no mueve datos. Solo se mueven las
  sumas diarias por reloj, una fila por reloj y día, para juntarlas por
  usuario.
\end{itemize}

\pregunta{¿Justifica hoy el caso un procesamiento distribuido?} No todavía.
Los datos actuales pesan 296 MB (E1) y el pipeline tarda 16.4 s en bronze y silver
en una máquina. Con todos los usuarios de Food.com, el lote diario sigue
cabiendo en una máquina si se procesa por archivo horario. Spark se justifica en
el escenario x100: un solo archivo horario necesitaría unos 378 GB de memoria
y volver a procesar los 3 años tardaría 309 días en una máquina.

\anexo{1}{Bibliografía}

\textbf{Material del curso} (Sebastián Hagedorn, INF3862 Plataformas Big Data
para Data Science, 2026):
\begin{enumerate}
  \item \emph{Clase 01: Almacenamiento de Datos: Fundamentos y Aplicación},
  págs. 26, 27 y 40: almacenamiento por columnas, compresión y
  particionamiento (5.3).
  \item \emph{Clase 02: ETL y Orquestación: De los datos al pipeline}, págs.
  24 y 26: idempotencia y ELT (5.1).
  \item \emph{Clase 04: Procesamiento con Spark}, págs. 15 y 21: shuffle al
  agrupar y broadcast en joins (5.4).
\end{enumerate}
\endgroup

\end{document}
